\begin{dashboardbox}[先说结论]
\begin{houseviewbox}[当前判断]
AI长剧现在是一个\textbf{值得盯住的9月候选题材}，还称不上主线。D+1午间，芒果超媒和多只主板AI内容扩散股继续涨停，说明市场在交易这条消息；新增报道也明确提到芒果TV出品和制作该项目。真正缺的仍是主板项目直接关系股，所以主板龙头战法的标准答案仍为空。
\end{houseviewbox}
\begin{sourcequalitybox}[你现在该怎么理解]
有创业板权限：芒果超媒是情绪温度计，用来观察题材强弱。无创业板权限：先等主板出现项目直接关系股；掌阅科技、博纳影业和中广天择只反映广义AI内容扩散，不能接替项目龙头。\sourcenote{S04；S05；S09。}
\end{sourcequalitybox}
\end{dashboardbox}

\begin{exhibitbox}[展品：一句话交易地图]
\begin{tabularx}{\textwidth}{L{3.1cm}X}
\toprule
你看到的现象 & 该怎么处理 \\
\midrule
只有芒果超媒强、主板跟风冲高 & 当作创业板情绪脉冲；主板账户不找“平替”。 \\
主板出现项目直接关系股，首板后走出二板 & 把它列为高度龙头候选，进入D+1/D+2复核。 \\
主板直接关系股连续两日放量、价格收得住 & 把它列为容量候选，和高度股二选一观察。 \\
题材次日掉队、主板断板或催化未续上 & 题材假设结束，删除观察标签。 \\
\bottomrule
\end{tabularx}
\end{exhibitbox}

这份预案只做一件事：每天收盘后判断“明天还值不值得继续看”。它不预设具体价格、仓位或收益目标。
